% ── Methods: model and training ──────────────────────────────────────────────
% Fragment to \input{} into the paper. Needs amsmath, amssymb, booktabs, array.
% Every value below is read from configs/*.yaml and src/ (see the % refs).
% Lines marked CHECK need a value or a decision that is not in this repository.

\subsection{Habitat targets}

Each specimen image $i$ (one of three views: back, belly, side) inherits the
habitat profile of its species, given as the fraction of use of $C=10$ habitat
classes derived from the IUCN level-1 habitat map of Jung et al.\ (2020):
forest, savanna, shrubland, grassland, desert, marine intertidal, artificial
terrestrial, marine neritic/oceanic, wetlands and coastal/artificial aquatic,
and rocky areas/caves.
% CHECK: describe how the per-species fractions were computed (eBird-weighted, data/raw/PHA_Jung_lvl1_fraction_v0.csv)
Fractions are renormalised to sum to one,
\begin{equation}
  \mathbf{y}_i \in \Delta^{C-1}, \qquad y_{ic} \ge 0,\quad \sum_{c=1}^{C} y_{ic} = 1,
\end{equation}
and the binary support (habitat used at all) is $s_{ic} = \mathbb{1}[y_{ic} > 0]$.
% src/datasets/folds.py: load_fold

Images are split into three folds at the family level: no family appears in both
the training and validation split of a fold, and each image is in exactly one
validation split. One model is trained per fold.
% data/cross/cross_version_1_FAM, checked by folds.check_partition
% CHECK: insert final image/species/family counts per fold (CSV rows are filtered to images present on disk)

\subsection{Bags of crops}

The model sees a specimen as a bag of $N = 16$ square crops taken from the
segmented (background-removed) image: $8$ crops of side $P = 128$ px, upsampled
to $224$ px with bicubic interpolation, and $8$ crops of side $P = 224$ px.
Crop positions are drawn from a precomputed table of admissible top-left
corners per specimen and crop size.
% CHECK: state the admissibility criterion of the .cand_{128,224} tables (built outside this repo)
To spread crops over the plumage, the admissible positions are grouped into grid
cells of side $P/2$; cells are visited in random order and one position is drawn
uniformly in each, so that crops of a given size start at least $P/2$ apart.
When a specimen has fewer cells than crops, the remaining positions are drawn
uniformly among all admissible positions.
% src/sampling.py

Every bag goes through the same fixed preprocessing chain:
(i) a median filter ($15 \times 15$) applied to the image before cropping;
(ii) a vertical flip with probability $0.5$;
(iii) cropping and resizing;
(iv) per-channel ImageNet normalisation.
Ablations (below) are inserted at fixed points of this chain
(Table~\ref{tab:ablations}).
% src/chain.py: build_bag, MEDIAN_BLUR_K = 15

\subsection{Online ablations}

During training, each bag is assigned a \emph{condition}: a set
$\mathcal{A} \subseteq \{A_1,\dots,A_7\}$ of image manipulations, drawn anew for
every image at every epoch. The number of ablations is drawn first,
\begin{equation}
  K \sim \mathrm{Categorical}\big(p_0, p_1, p_2, p_3\big) = (0.20,\ 0.40,\ 0.30,\ 0.10),
\end{equation}
then $\mathcal{A}$ is drawn uniformly among the admissible sets of size $K$.
A set is admissible if it contains at most one of the four colour ablations
$\{A_1, A_2, A_6, A_7\}$, which act on the same CIELAB channels and would
otherwise overwrite one another. This gives $1 + 7 + 15 + 13 = 36$ admissible
conditions. The marginal probability that a given ablation is active in a bag
is $0.140$ for each colour ablation and $0.246$ for each of $A_3$, $A_4$, $A_5$;
the unaltered reference condition ($\mathcal{A} = \varnothing$) is seen in
$20\%$ of bags.
% src/ablations/combinations.py: K_PROBABILITIES, marginal_inclusion_probabilities

All random draws for a bag (condition, ablation parameters, crop positions,
flip) come from a generator seeded by a hash of
(run seed, epoch, image index, image name), so training is reproducible
regardless of data-loader scheduling.
% src/rng.py, src/datasets/bag_dataset.py

\begin{table}[t]
\centering
\small
\begin{tabular}{@{}l >{\raggedright\arraybackslash}p{5.6cm} l >{\raggedright\arraybackslash}p{5.0cm}@{}}
\toprule
Code & Manipulation & Level & Parameters (drawn per) \\
\midrule
$A_1$ & Achromatic: $a^* = b^* = 0$, $L^*$ kept & image & -- \\
$A_2$ & $L^*$ matched to the population mean, s.d.\ and histogram; $a^*, b^*$ kept & image & -- \\
$A_3$ & Random permutation of $h \times h$ blocks & crop & $h \in \{1,2,4,7,8,14,16,28,32,56\}$ (crop) \\
$A_4$ & Directional box blur & image & $\theta \in \{0^\circ, 90^\circ\}$, $\ell \in [17, 101]$ (bag) \\
$A_5$ & Amplitude-spectrum slope set to a reference & crop & $\alpha^\star = \bar\alpha_P + \mathcal{U}(-0.1, 0.1)$ (crop) \\
$A_6$ & Chroma of one colour category removed, $L^*$ kept & image & category present on specimen (bag) \\
$A_7$ & Isoluminant: $L^* = 50$, $a^*, b^*$ kept & image & -- \\
\bottomrule
\end{tabular}
\caption{Ablations drawn online during training. CIELAB conversions use the
D65 white point; colour ablations act on the foreground only.}
\label{tab:ablations}
\end{table}

\paragraph{Colour ablations.}
With $(L^*, a^*, b^*)$ the CIELAB coordinates of a foreground pixel,
$A_1$ maps it to $(L^*, 0, 0)$ and $A_7$ to $(50, a^*, b^*)$, followed by
conversion back to sRGB with gamut clipping.
$A_2$ standardises lightness across the dataset in the manner of the SHINE
toolbox (Willenbockel et al., 2010): each specimen's $L^*$ is first rescaled to
the population mean and standard deviation, then exactly histogram-matched to a
population target histogram (each specimen weighted equally), leaving $a^*$ and
$b^*$ unchanged.
% CHECK: the evaluated runs used the cluster A2 cache, which is greyscale (memory: cluster-dataset-differs). Align this sentence with that version.
$A_6$ classifies every foreground pixel into the 12 avian colour categories of
Delhey et al.\ (2023), via the CN11 colour names of van de Weijer et al.\ (2009);
one of the chromatic categories present on the specimen is drawn uniformly and
its pixels are set to $(L^*, 0, 0)$.
% src/ablations/a1_*.py, a2_*.py, a6_*.py, a7_*.py; A1, A2, A7 and the A6 label maps are precomputed

\paragraph{Spatial ablations.}
$A_3$ splits a $224$-px crop into $(224/h)^2$ blocks of side $h$ and shuffles
them with a uniform random permutation, preserving the colour histogram.
$A_4$ convolves the image with a $1 \times \ell$ ($\theta = 0^\circ$) or
$\ell \times 1$ ($\theta = 90^\circ$) box kernel, with $\ell$ drawn
log-uniformly on $[17, 101]$ and rounded to an odd integer; it is applied
before cropping so that the blur sees real neighbourhoods.
$A_5$ regrades the amplitude spectrum of each crop. Let $\hat g(\mathbf{k})$ be
the 2-D Fourier transform of the green channel and $\bar A(r)$ its radially
averaged amplitude at integer radius $r = \lfloor \lVert\mathbf{k}\rVert \rceil$.
The target profile is the power law
\begin{equation}
  A^\star(r) = \bar A(r_0)\,\left(\frac{r}{r_0}\right)^{-\alpha^\star}, \qquad r_0 = 2,
\end{equation}
and every channel is multiplied in Fourier space by the real gain
$G(\mathbf{k}) = A^\star(r) / \bar A(r)$, with $G = 1$ at the DC term. Phases
and mean brightness are unchanged. The target slope $\alpha^\star$ is centred on
the mean slope measured on reference crops of the same source size
($\bar\alpha_{224} = 1.49$, $\bar\alpha_{128} = 2.19$, fitted by least squares in
log--log space over $r \in [2, 9]$).
% src/ablations/a3_*.py, a4_*.py, a5_*.py

\paragraph{Sham control.}
For each fold we also train a sham model with the identical pipeline, seed and
hyperparameters, except that every bag is in the reference condition
($\mathcal{A} = \varnothing$). The median filter and flip are kept.
% configs/train_sham_fold*.yaml

\subsection{Model}

\paragraph{Backbone.}
Each crop $x_{ij}$ is embedded by a DINOv3 ViT-S/16 (12 blocks, width
$d = 384$; Siméoni et al., 2025); we use its normalised class token,
$\mathbf{h}_{ij} = f_\theta(x_{ij}) \in \mathbb{R}^{d}$.
The pretrained weights $W_0$ are frozen and adapted with LoRA (Hu et al., 2022)
on the attention input ($\mathrm{qkv}$) and output ($\mathrm{proj}$) projections
of every block:
\begin{equation}
  W = W_0 + \frac{\alpha}{r}\, B A, \qquad A \in \mathbb{R}^{r \times d_\text{in}},\ B \in \mathbb{R}^{d_\text{out} \times r},
\end{equation}
with rank $r = 8$, $\alpha = 18$ and dropout $0.05$ on the LoRA input.
% configs/model.yaml, src/model/backbone.py

\paragraph{Gated attention pooling.}
The $N$ crop embeddings of a bag are pooled with the gated attention of
Ilse et al.\ (2018):
\begin{align}
  e_{ij} &= \mathbf{w}^\top \big( \tanh(\mathbf{V}\mathbf{h}_{ij} + \mathbf{b}_V) \odot \sigma(\mathbf{U}\mathbf{h}_{ij} + \mathbf{b}_U) \big) + b_w, \\
  a_{ij} &= \frac{\exp e_{ij}}{\sum_{k=1}^{N} \exp e_{ik}}, \qquad
  \mathbf{z}_i = \sum_{j=1}^{N} a_{ij}\, \mathbf{h}_{ij},
\end{align}
with $\mathbf{V}, \mathbf{U} \in \mathbb{R}^{96 \times 384}$ ($d/4$ hidden units).
% src/model/mil.py

\paragraph{Heads.}
Two heads read the bag embedding, each a two-layer MLP
($384 \to 192$, GELU, dropout $0.1$, $192 \to C$):
a regression head giving the predicted habitat distribution
$\hat{\mathbf{p}}_i = \mathrm{softmax}(\mathbf{r}_i)$, and a classification head
giving per-habitat presence probabilities $\hat q_{ic} = \sigma(u_{ic})$.
% src/model/heads.py

\subsection{Loss}

The regression head is trained with the Kullback--Leibler divergence from the
target distribution, and the classification head with a class-weighted binary
cross-entropy on the support:
\begin{align}
  \mathcal{L}_\text{KL} &= \frac{1}{|\mathcal{B}|} \sum_{i \in \mathcal{B}} \sum_{c=1}^{C} y_{ic} \big( \log y_{ic} - \log \hat p_{ic} \big), \\
  \mathcal{L}_\text{BCE} &= \frac{1}{|\mathcal{B}|\,C} \sum_{i \in \mathcal{B}} \sum_{c=1}^{C} w_c \Big[ -s_{ic} \log \hat q_{ic} - (1 - s_{ic}) \log (1 - \hat q_{ic}) \Big],
\end{align}
with $0 \log 0 = 0$. The class weights are inverse support frequencies on the
training split of the fold, normalised to mean one:
\begin{equation}
  f_c = \frac{1}{n} \sum_{i=1}^{n} s_{ic}, \qquad
  w_c = \frac{1/f_c}{\frac{1}{C} \sum_{c'} 1/f_{c'}}.
\end{equation}
The total loss is
$\mathcal{L} = \lambda_\text{KL}\, \mathcal{L}_\text{KL} + \lambda_\text{BCE}\, \mathcal{L}_\text{BCE}$,
with $\lambda_\text{KL} + \lambda_\text{BCE} = 1$.
% src/model/losses.py, src/datasets/folds.py: bce_class_weights

\paragraph{Adaptive loss weights.}
The weights start at $0.5/0.5$ and are updated once per epoch so that the term
whose loss is decreasing more slowly receives more weight, in the spirit of
dynamic weight averaging (Liu et al., 2019). Let $\ell_k^{(e)}$ be the mean of
term $k \in \{\text{KL}, \text{BCE}\}$ over epoch $e$. We track an exponential
moving average and its relative change over $\Delta = 5$ epochs,
\begin{equation}
  m_k^{(e)} = \beta\, m_k^{(e-1)} + (1 - \beta)\, \ell_k^{(e)}, \quad m_k^{(1)} = \ell_k^{(1)}, \qquad
  \rho_k^{(e)} = \frac{m_k^{(e)}}{m_k^{(e-\Delta)}},
\end{equation}
with $\beta = 0.9$, and set the weights for epoch $e + 1$ to
\begin{equation}
  \lambda_\text{KL}^{(e+1)} = \mathrm{clip}\!\left( \sigma\!\left( \frac{\rho_\text{KL}^{(e)} - \rho_\text{BCE}^{(e)}}{T} \right),\ 0.2,\ 0.8 \right),
  \qquad \lambda_\text{BCE}^{(e+1)} = 1 - \lambda_\text{KL}^{(e+1)},
\end{equation}
with temperature $T = 0.1$ (the sigmoid is the two-term softmax of $\rho_k / T$).
Weights stay at $0.5/0.5$ until $\Delta + 1$ epochs of history are available,
i.e.\ for epochs $1$--$6$.
% src/training/loss_balancing.py, configs/optim.yaml

\subsection{Optimisation}

Only the LoRA matrices, the attention pooling and the heads are trained. We use
AdamW (Loshchilov \& Hutter, 2019) with learning rate $\eta_0 = 10^{-4}$ and
decoupled weight decay $0.01$ on weight matrices; biases and all attention-pool
parameters are not decayed (decaying the pool drove the attention towards the
uniform $1/N$). The learning rate follows a cosine schedule over
$E_\text{max} = 100$ epochs, updated per epoch,
\begin{equation}
  \eta_e = \frac{\eta_0}{2} \left( 1 + \cos \frac{\pi (e - 1)}{E_\text{max}} \right),
\end{equation}
whether or not training stops early. An epoch is one pass over the training
images, each forming one bag with a freshly drawn condition. Each optimiser step
uses $128$ bags ($2048$ crops), accumulated over $8$ micro-batches of $16$ bags;
forward passes run in bfloat16 mixed precision.
% src/training/run.py, configs/optim.yaml

\paragraph{Stopping and model selection.}
Selection uses only the training loss. At each epoch we compute the training
loss with fixed weights,
$\bar{\mathcal{L}}^{(e)} = 0.5\, \ell_\text{KL}^{(e)} + 0.5\, \ell_\text{BCE}^{(e)}$,
so that it is not affected by the adaptive reweighting, and average it over the
last $5$ epochs. The checkpoint with the lowest smoothed value is kept, and
training stops after $20$ epochs without improvement or at epoch $100$. The loss
on reference-condition bags of the validation families is logged for monitoring
but is never used for selection or stopping.
% src/training/run.py: fit

All runs (three folds, main and sham) use the same seed. For each run we store
the resolved configuration, the code commit, library versions and the class
weights.
% seed 20260911, configs/train_*.yaml
